% GENERATED FILE. Edit canonical lessons, then run npm run generate:books.
% canonical-source-hash: fnv1a64:afb291511e841fa9
% canonical-lessons: KA-C408-kendrikarisu, KA-C408-adhyayana, KA-C408-aicchika, KA-C408-pathyakrama, KA-C408-punaravartane

\chapter{Studying Well}
\label{ch:ka-studying-well}
\hlchaptermodality{408}

\begin{chapteropening}
\textbf{By the end of this chapter:} \emph{I can say to focus, a study, optional, a syllabus and revision.}

Complete the last lesson of chapter 408: I can say to focus, a study, optional, a syllabus and revision.
\end{chapteropening}

\section[\texorpdfstring{\kn{ಕೇಂದ್ರೀಕರಿಸು} (kēndrīkarisu)}{ಕೇಂದ್ರೀಕರಿಸು (kēndrīkarisu)}]{\kn{ಕೇಂದ್ರೀಕರಿಸು} (kēndrīkarisu) — to focus}

\label{lesson:KA-C408-kendrikarisu}



\begin{quote}
Before the new one: say the Kannada for to analyse, then the Kannada for a survey.
\end{quote}

\subsection*{You'll want to know: \kn{ಕೇಂದ್ರೀಕರಿಸು}}
\textbf{\kn{ಕೇಂದ್ರೀಕರಿಸು}} — \emph{kēndrīkarisu} — \textquotedblleft{}to focus\textquotedblright{}.

\begin{quote}
\kn{ಪರೀಕ್ಷೆಗೆ} \kn{ಮೊದಲು} \kn{ನಾನು} \kn{ಗಣಿತದ} \kn{ಮೇಲೆ} \kn{ಗಮನ} \kn{ಕೇಂದ್ರೀಕರಿಸಿದೆ}.
\end{quote}

\begin{quote}
\emph{parīkṣege modalu nānu gaṇitada mēle gamana kēndrīkariside.}
\end{quote}

\begin{quote}
Before the exam I focused on mathematics.
\end{quote}

Past (I …): \kn{ಕೇಂದ್ರೀಕರಿಸಿದೆ}, \emph{kēndrīkariside}.

What you focus on takes \kn{ಮೇಲೆ}.

\begin{cousinweb}
\textbf{\kn{ಕೇಂದ್ರೀಕರಿಸು}} (\emph{kēndrīkarisu}) is built on Sanskrit \emph{kendra}, \textquotedblleft{}the centre of a circle\textquotedblright{} (Monier-Williams), itself taken from Greek \emph{kentron}, with -\emph{īkarisu}, \textquotedblleft{}to make into\textquotedblright{}.
\end{cousinweb}

\begin{grammarlens}[title={What you've built}]
One more word for giving reasons and opinions.
\end{grammarlens}

\subsection*{Your turn}
\begin{itemize}
\raggedright
  \item \emph{Say it:} \emph{kēndrīkarisu}
  \item \emph{Say it:} \emph{kēndrīkarisu}, once more
  \item \emph{Say it:} \emph{parīkṣege modalu nānu gaṇitada mēle gamana kēndrīkariside.} — Before the exam I focused on mathematics.
  \item \emph{Recall:} say the Kannada for to analyse, then the Kannada for a survey, then say \emph{kēndrīkarisu} again
\end{itemize}

\begin{tcolorbox}[breakable,colback=teal!4,colframe=teal!35!black,title={Before you move on}]
What does \kn{ಕೇಂದ್ರೀಕರಿಸು} mean? (\textquotedblleft{}To focus\textquotedblright{}.) Say it once more.
\end{tcolorbox}
\section[\texorpdfstring{\kn{ಅಧ್ಯಯನ} (adhyayana)}{ಅಧ್ಯಯನ (adhyayana)}]{\kn{ಅಧ್ಯಯನ} (adhyayana) — a study; studying}

\label{lesson:KA-C408-adhyayana}



\begin{quote}
Before the new one: say the Kannada for a survey, then the Kannada for to focus.
\end{quote}

\subsection*{You'll want to know: \kn{ಅಧ್ಯಯನ}}
\textbf{\kn{ಅಧ್ಯಯನ}} — \emph{adhyayana} — \textquotedblleft{}a study; studying\textquotedblright{}.

\begin{quote}
\kn{ಅವನು} \kn{ಹಳ್ಳಿಯ} \kn{ನೀರಿನ} \kn{ಬಗ್ಗೆ} \kn{ಅಧ್ಯಯನ} \kn{ಮಾಡಿದನು}.
\end{quote}

\begin{quote}
\emph{avanu haḷḷiya nīrina bagge adhyayana māḍidanu.}
\end{quote}

\begin{quote}
He made a study of the village's water.
\end{quote}

\kn{ಅಧ್ಯಯನ} \kn{ಮಾಡು}, to study (a subject in depth).

\begin{cousinweb}
\textbf{\kn{ಅಧ್ಯಯನ}} (\emph{adhyayana}) is borrowed from Sanskrit \emph{adhyayana}, \textquotedblleft{}reading, studying\textquotedblright{} (Monier-Williams).
\end{cousinweb}

\begin{grammarlens}[title={What you've built}]
One more word for giving reasons and opinions.
\end{grammarlens}

\subsection*{Your turn}
\begin{itemize}
\raggedright
  \item \emph{Say it:} \emph{adhyayana}
  \item \emph{Say it:} \emph{adhyayana}, once more
  \item \emph{Say it:} \emph{avanu haḷḷiya nīrina bagge adhyayana māḍidanu.} — He made a study of the village's water.
  \item \emph{Recall:} say the Kannada for a survey, then the Kannada for to focus, then say \emph{adhyayana} again
\end{itemize}

\begin{tcolorbox}[breakable,colback=teal!4,colframe=teal!35!black,title={Before you move on}]
What does \kn{ಅಧ್ಯಯನ} mean? (\textquotedblleft{}A study; studying\textquotedblright{}.) Say it once more.
\end{tcolorbox}
\section[\texorpdfstring{\kn{ಐಚ್ಛಿಕ} (aicchika)}{ಐಚ್ಛಿಕ (aicchika)}]{\kn{ಐಚ್ಛಿಕ} (aicchika) — optional; elective}

\label{lesson:KA-C408-aicchika}



\begin{quote}
Before the new one: say the Kannada for to focus, then the Kannada for a study.
\end{quote}

\subsection*{You'll want to know: \kn{ಐಚ್ಛಿಕ}}
\textbf{\kn{ಐಚ್ಛಿಕ}} — \emph{aicchika} — \textquotedblleft{}optional; elective\textquotedblright{}.

\begin{quote}
\kn{ಶಾಲೆಯಲ್ಲಿ} \kn{ಕನ್ನಡ} \kn{ಕಡ್ಡಾಯ}, \kn{ಸಂಗೀತ} \kn{ಐಚ್ಛಿಕ}.
\end{quote}

\begin{quote}
\emph{śāleyalli kannaḍa kaḍḍāya, saṅgīta aicchika.}
\end{quote}

\begin{quote}
At school Kannada is compulsory and music is optional.
\end{quote}

Its opposite is \kn{ಕಡ್ಡಾಯ}, compulsory. \kn{ಐಚ್ಛಿಕ} \kn{ವಿಷಯ}, an elective subject.

\begin{cousinweb}
\textbf{\kn{ಐಚ್ಛಿಕ}} (\emph{aicchika}) is borrowed from Sanskrit \emph{aicchika}, \textquotedblleft{}optional, arbitrary, at will\textquotedblright{} (Monier-Williams), an adjective formed from \emph{icchā}, \textquotedblleft{}a wish\textquotedblright{}.
\end{cousinweb}

\begin{grammarlens}[title={What you've built}]
One more word for giving reasons and opinions.
\end{grammarlens}

\subsection*{Your turn}
\begin{itemize}
\raggedright
  \item \emph{Say it:} \emph{aicchika}
  \item \emph{Say it:} \emph{aicchika}, once more
  \item \emph{Say it:} \emph{śāleyalli kannaḍa kaḍḍāya, saṅgīta aicchika.} — At school Kannada is compulsory and music is optional.
  \item \emph{Recall:} say the Kannada for to focus, then the Kannada for a study, then say \emph{aicchika} again
\end{itemize}

\begin{tcolorbox}[breakable,colback=teal!4,colframe=teal!35!black,title={Before you move on}]
What does \kn{ಐಚ್ಛಿಕ} mean? (\textquotedblleft{}Optional; elective\textquotedblright{}.) Say it once more.
\end{tcolorbox}
\section[\texorpdfstring{\kn{ಪಠ್ಯಕ್ರಮ} (paṭhyakrama)}{ಪಠ್ಯಕ್ರಮ (paṭhyakrama)}]{\kn{ಪಠ್ಯಕ್ರಮ} (paṭhyakrama) — a syllabus}

\label{lesson:KA-C408-pathyakrama}



\begin{quote}
Before the new one: say the Kannada for a study, then the Kannada for optional.
\end{quote}

\subsection*{You'll want to know: \kn{ಪಠ್ಯಕ್ರಮ}}
\textbf{\kn{ಪಠ್ಯಕ್ರಮ}} — \emph{paṭhyakrama} — \textquotedblleft{}a syllabus\textquotedblright{}.

\begin{quote}
\kn{ಈ} \kn{ವರ್ಷ} \kn{ಪಠ್ಯಕ್ರಮ} \kn{ಬದಲಾಯಿತು}.
\end{quote}

\begin{quote}
\emph{ī varṣa paṭhyakrama badalāyitu.}
\end{quote}

\begin{quote}
The syllabus changed this year.
\end{quote}

\kn{ಪಠ್ಯ}, a text to be read; \kn{ಪಠ್ಯಪುಸ್ತಕ}, a textbook.

\begin{cousinweb}
\textbf{\kn{ಪಠ್ಯಕ್ರಮ}} (\emph{paṭhyakrama}) is a modern word built from Sanskrit \emph{paṭhya}, \textquotedblleft{}to be read\textquotedblright{}, from \emph{paṭh}, \textquotedblleft{}to read\textquotedblright{}, and \emph{krama}, \textquotedblleft{}a step, an order\textquotedblright{} (CDIAL 3577).
\end{cousinweb}

\begin{grammarlens}[title={What you've built}]
One more word for giving reasons and opinions.
\end{grammarlens}

\subsection*{Your turn}
\begin{itemize}
\raggedright
  \item \emph{Say it:} \emph{paṭhyakrama}
  \item \emph{Say it:} \emph{paṭhyakrama}, once more
  \item \emph{Say it:} \emph{ī varṣa paṭhyakrama badalāyitu.} — The syllabus changed this year.
  \item \emph{Recall:} say the Kannada for a study, then the Kannada for optional, then say \emph{paṭhyakrama} again
\end{itemize}

\begin{tcolorbox}[breakable,colback=teal!4,colframe=teal!35!black,title={Before you move on}]
What does \kn{ಪಠ್ಯಕ್ರಮ} mean? (\textquotedblleft{}A syllabus\textquotedblright{}.) Say it once more.
\end{tcolorbox}
\section[\texorpdfstring{\kn{ಪುನರಾವರ್ತನೆ} (punarāvartane)}{ಪುನರಾವರ್ತನೆ (punarāvartane)}]{\kn{ಪುನರಾವರ್ತನೆ} (punarāvartane) — revision; repetition}

\label{lesson:KA-C408-punaravartane}



\begin{quote}
Before the new one: say the Kannada for optional, then the Kannada for a syllabus.
\end{quote}

\subsection*{You'll want to know: \kn{ಪುನರಾವರ್ತನೆ}}
\textbf{\kn{ಪುನರಾವರ್ತನೆ}} — \emph{punarāvartane} — \textquotedblleft{}revision; repetition\textquotedblright{}.

\begin{quote}
\kn{ಪರೀಕ್ಷೆಗೆ} \kn{ಮೊದಲು} \kn{ಪುನರಾವರ್ತನೆ} \kn{ಮುಖ್ಯ}.
\end{quote}

\begin{quote}
\emph{parīkṣege modalu punarāvartane mukhya.}
\end{quote}

\begin{quote}
Revision is important before an exam.
\end{quote}

\kn{ಪುನರ್} (again) + \kn{ಆವರ್ತನೆ} (turning round).

\begin{cousinweb}
\textbf{\kn{ಪುನರಾವರ್ತನೆ}} (\emph{punarāvartane}) is built from Sanskrit \emph{punar}, \textquotedblleft{}again\textquotedblright{}, and \emph{āvartana}, \textquotedblleft{}turning round, revolving\textquotedblright{} (Monier-Williams).
\end{cousinweb}

\begin{grammarlens}[title={What you've built}]
One more word for giving reasons and opinions.
\end{grammarlens}

\subsection*{Your turn}
\begin{itemize}
\raggedright
  \item \emph{Say it:} \emph{punarāvartane}
  \item \emph{Say it:} \emph{punarāvartane}, once more
  \item \emph{Say it:} \emph{parīkṣege modalu punarāvartane mukhya.} — Revision is important before an exam.
  \item \emph{Recall:} say the Kannada for optional, then the Kannada for a syllabus, then say \emph{punarāvartane} again
\end{itemize}

\begin{tcolorbox}[breakable,colback=teal!4,colframe=teal!35!black,title={Before you move on}]
What does \kn{ಪುನರಾವರ್ತನೆ} mean? (\textquotedblleft{}Revision; repetition\textquotedblright{}.) Say it once more.
\end{tcolorbox}
